\documentclass[12pt]{article}
\usepackage[T1]{fontenc}
\usepackage[utf8]{inputenc}
\usepackage{lmodern}
\usepackage{textcomp}
\usepackage{enumitem}
\usepackage{geometry}
\geometry{margin=1in}
\begin{document}
\begin{center}
Answer Key - Philosophy - Undergraduate Level Assessment 17 \
\small (Answers are ordered exactly as in the question paper.)
\end{center}

\section*{Multiple Choice}
\begin{enumerate}[label=\arabic*.]
\item \textbf{Answer:} A
\\ \textit{Options:} A) Ashoka; B) Nelson Mandela; C) Martin Luther King Jr.; D) Ngo Dinh Diem; E) Adolf Hitler
\\ \textit{Note:} Ashoka is recognized for promoting the principles of non-violence and compassion through his edicts and policies after his conversion to Buddhism, which fundamentally emphasized the importance of dharma as a way to achieve peace and coexistence.
\item \textbf{Answer:} D
\\ \textit{Options:} A) Sound; B) Non-living; C) Non-matter; D) Non-soul; E) Consciousness
\\ \textit{Note:} The term ajiva in Jaina traditions refers to all non-soul entities or matter, distinguishing them from jiva, which signifies living souls, thereby highlighting the fundamental philosophical distinction between sentient and insentient beings.
\item \textbf{Answer:} A
\\ \textit{Options:} A) tu quoque; B) appeal to ignorance; C) argumentum ad populum; D) false cause; E) appeal to indignation
\\ \textit{Note:} The correct answer is tu quoque because the sarcastic response dismisses the argument by attacking the speaker's credibility or intentions rather than addressing the argument itself, which is characteristic of this fallacy.
\item \textbf{Answer:} C
\\ \textit{Options:} A) not adhering to societal norms.; B) being dishonest in any situation.; C) failure to perform one's covenant.; D) treating another person as a mere means.; E) manipulating others for personal gain.
\\ \textit{Note:} Hobbes defines injustice as the failure to perform one's covenant because he believes that social contracts are the foundation of moral and legal obligations, and breaking these covenants undermines the trust and order necessary for a functioning society.
\item \textbf{Answer:} A
\\ \textit{Options:} A) determine which pleasure most experienced people prefer; B) consult science; C) consult religious leaders; D) determine which one is objectively most pleasurable; E) measure the intensity of each pleasure
\\ \textit{Note:} This answer is correct because Mill emphasizes that the quality of pleasures can be assessed by considering the preferences of those who have experienced both types, as they are best positioned to judge which brings greater satisfaction.
\end{enumerate}

\section*{True / False}
\begin{enumerate}[label=\arabic*.]
\setcounter{enumi}{5}
\item \textbf{Answer:} False
\\ \textit{Options:} A) True; B) False
\\ \textit{Note:} Nagel argues that while absolutism emphasizes the moral prohibition against murder, it allows for exceptions in extreme cases where other moral considerations may take precedence.
\item \textbf{Answer:} False
\\ \textit{Options:} A) True; B) False
\\ \textit{Note:} Augustine challenges Academic skepticism by arguing that knowledge of certain truths, particularly those related to God and moral principles, is attainable, thus refuting the claim that such skepticism is unchallengeable or beyond disproof.
\item \textbf{Answer:} True
\\ \textit{Options:} A) True; B) False
\\ \textit{Note:} Socrates believed that the pursuit of self-knowledge and critical reflection is essential for the well-being of the soul, and living an unexamined life leads to ignorance and moral failings, thereby causing harm to one's inner self.
\item \textbf{Answer:} True
\\ \textit{Options:} A) True; B) False
\\ \textit{Note:} According to Aristotle, virtue and vice are the results of our choices and actions, as he believes that individuals have the capacity to cultivate moral character through deliberate practice and decision-making.
\item \textbf{Answer:} True
\\ \textit{Options:} A) True; B) False
\\ \textit{Note:} Hume argues that moral judgments arise from human sentiments and emotions shaped by cultural norms rather than from rational principles, emphasizing the role of social context in determining morality.
\end{enumerate}

\section*{Long Form Response}
\begin{enumerate}[label=\arabic*.]
\setcounter{enumi}{10}
\item \textbf{Answer:} Ziran, often translated as "naturalness" or "spontaneity," is a central concept in Zhuangzi's philosophy that emphasizes living in accordance with one's true nature and the natural order of the world. This notion has profound implications for compassion in human relationships, as it encourages individuals to act authentically and empathetically, allowing for genuine connections without the constraints of societal expectations. By fostering a sense of interconnectedness and understanding, ziran promotes a compassionate approach to interpersonal dynamics, suggesting that true kindness arises when individuals are free to express their inherent nature. In philosophical discourse, ziran challenges rigid moral frameworks, advocating instead for a fluid and adaptive understanding of ethics that aligns with the complexities of human experience. Thus, ziran not only enriches our understanding of compassion but also invites deeper contemplation on the nature of morality and human existence.
\\ \textit{Note:} The answer correctly identifies ziran as a fundamental principle in Zhuangzi's philosophy that promotes authenticity and empathy in human relationships, thereby enhancing compassion and fostering genuine connections aligned with the natural order.
\item \textbf{Answer:} Ikebana, the Japanese art of flower arranging, holds significant philosophical value as it embodies principles of simplicity, balance, and harmony with nature. Its practice encourages a mindful engagement with the natural world, reflecting broader themes in aesthetics that prioritize the ephemeral beauty of life and the interconnectedness of all things. Each arrangement is a meditative process that underscores the importance of presence and intentionality, inviting practitioners and observers alike to appreciate the fleeting nature of beauty. Moreover, Ikebana challenges Western conventions of beauty that often emphasize abundance and perfection, instead celebrating asymmetry and minimalism. Ultimately, Ikebana serves as a philosophical lens through which we can explore deeper existential questions about the nature of beauty, impermanence, and our relationship with the environment.
\\ \textit{Note:} correct because it effectively articulates how Ikebana embodies philosophical principles such as simplicity and balance, highlights the importance of mindfulness and intentionality, and underscores broader themes in aesthetics concerning the transient beauty of nature and the interconnectedness of all living things.
\end{enumerate}

\vfill
\noindent\textit{End of Answer Key}
\end{document}